\documentclass[12pt]{book}
\usepackage{amsmath}
\usepackage{amssymb}
\usepackage{geometry}
\geometry{margin=1in}

\title{Volume 03: Atomic Physics (Summary Only; Full Canon in ATOMICUS) \\ Book 01: Atomic Structure Overview}
\author{James C. Harvey}
\date{December 2025}

\begin{document}
\maketitle
\tableofcontents

\chapter{Atomic Structure in SDT: Structural Summary}

\section*{Abstract}
This book provides a structural summary of atomic physics within SDT. It does not replicate the full
element-by-element canon, which resides in \texttt{SDT/ATOMICUS}. Instead, it establishes the geometric
grammar of atomic structure, the matrix--flux coupling that binds electrons, and the role of occlusion
and circulation in defining discrete shells.

\section{Scope and Canon Boundary}
Atomic physics is complete in \texttt{SDT/ATOMICUS}. This book defines the structural frame: how an atom
is represented in SDT and how discrete behavior follows from geometry and flux.

\section{Atomic Binding as Flux Balance}
At the Bohr scale, binding is the balance of flux pressure and boundary curvature. The effective coupling
is a geometric occlusion limit:
\begin{equation}
F_{\text{atomic}} = \frac{\kappa_1 \kappa_2}{4\pi K_{\text{bulk}} r^2},
\end{equation}
with $\kappa$ determined by boundary curvature and $r$ set by orbital geometry.

\section{Shells as Helical Path Closure}
Discrete shells are helical path closures where matrix flux and boundary circulation are phase-locked.
The orbital period is
\begin{equation}
T = \frac{2\pi r}{v},
\end{equation}
and stable shells occur when this period matches the boundary’s internal oscillation mode.

\section{Occlusion and Discrete States}
Occlusion is the state selector. The occlusion function
\begin{equation}
E(\mathbf{r},\hat{\mathbf{n}}) = \frac{\Omega_{\text{blocked}}(\mathbf{r},\hat{\mathbf{n}})}{4\pi}
\end{equation}
determines whether the interaction is in a Coulomb-like regime ($E \to 0$) or a screened regime ($E \to 1$).
Discrete states are stabilized by occlusion symmetry and circulation lock.

\section{Helical Wakes and Magnetic Structure}
The helical wake is the magnetic signature of circulation. The magnetic moment is proportional to
circulation and curvature:
\begin{equation}
\mu \propto \Gamma \kappa (1-\eta).
\end{equation}
This governs hyperfine interactions and the spectral structure addressed in the atomic canon.

\section{Numerical Template for Moments}
For atomic-scale boundaries, the moment is evaluated using the conversion constant $C_\mu$:
\begin{equation}
\mu = C_\mu \Gamma \kappa (1-\eta).
\end{equation}
All atomic evaluations in \texttt{SDT/ATOMICUS} use this template with geometry-derived $\Gamma$, $\kappa$,
and $\eta$.

\section{ATOMICUS Cross-Reference}
All element-specific geometry, isotopic structure, and atomic properties are defined in \texttt{SDT/ATOMICUS}.
This book is the structural frame; the canon is the full atlas.

\chapter{Geometric Grammar of the Atom}

\section*{Abstract}
Atomic structure in SDT is a geometric grammar: the nucleus is a compact boundary core, electrons are
toroidal vortex structures, and shells are flux-locked helices. This chapter provides the grammar without
repeating the element-by-element atlas.

\section{Electron as Toroidal Vortex}
Electrons are modeled as toroidal vortices with circulation $\Gamma$ and curvature $\kappa$:
\begin{equation}
\mu_e = C_\mu\, \Gamma \kappa (1-\eta),
\end{equation}
where $C_\mu$ is the geometric conversion from circulation to moment. The helical wake is the observable
magnetic imprint of this circulation.

\section{Nuclear Core as Compact Boundary}
The nucleus is a compact boundary region where occlusion is high and curvature is extreme. The binding
condition is the matching of flux capture to internal circulation:
\begin{equation}
\dot{E}_{\text{bind}} = P_{\text{CMB}} A_{\text{eff}} \Gamma \kappa (1-\eta).
\end{equation}

\section{Shell Quantization as Phase Lock}
Quantization is the closure of helical paths:
\begin{equation}
n \lambda_{\text{flux}} = 2\pi r,
\end{equation}
where $\lambda_{\text{flux}}$ is the effective flux wavelength imposed by internal oscillation. This
expression is a geometric phase condition rather than an external postulate.

\section{Quantization Derivative}
Differentiating the closure condition yields the discrete spacing rule:
\begin{equation}
\Delta r = \frac{\lambda_{\text{flux}}}{2\pi}\Delta n.
\end{equation}
This relation encodes shell spacing in terms of the internal flux wavelength.

\chapter{Isotopes, Spin, and Hyperfine Structure}

\section*{Abstract}
Isotopic variation is a change in boundary packing; hyperfine structure is a change in circulation
coupling. This chapter summarizes these mechanisms and points to the canonical atlas.

\section{Isotopic Variation}
Isotopes are changes in nuclear boundary topology that alter occlusion symmetry:
\begin{equation}
\Delta E_{\text{iso}} \propto \Delta \kappa \cdot \Gamma.
\end{equation}
The detailed structures and values are cataloged in \texttt{SDT/ATOMICUS}.

\section{Hyperfine Coupling}
Hyperfine splitting follows from coupling between electron circulation and nuclear boundary oscillation:
\begin{equation}
\Delta f_{\text{hf}} \propto \Gamma_e \kappa_n (1-\eta).
\end{equation}

\section{Scaling with Curvature}
Holding $\Gamma_e$ constant, the hyperfine shift scales linearly with nuclear curvature:
\begin{equation}
\frac{\Delta f_{\text{hf,2}}}{\Delta f_{\text{hf,1}}} = \frac{\kappa_{n,2}}{\kappa_{n,1}}.
\end{equation}

\chapter{Reading the Atomic Canon}

\section*{Abstract}
This chapter explains how to use the atomic canon. The treatise provides structure; the atlas provides
detail.

\section{Canonical Index}
Each element in \texttt{SDT/ATOMICUS} defines its geometry, circulation, occlusion, and derived moments.
The treatise does not replicate those values; it supplies the framework that makes them intelligible.

\section{Scope Boundary}
If a quantity is element-specific, it belongs in \texttt{SDT/ATOMICUS}. If it is structural and general,
it belongs in this volume. This boundary preserves clarity and prevents duplication.

\end{document}
